% Generated draft: requires mathematical and source review.
\documentclass[11pt]{article}
\usepackage[margin=25mm]{geometry}
\usepackage{fontspec}
\setmainfont{Latin Modern Roman}
\usepackage{amsmath,amssymb,mathtools,unicode-math}
\setmathfont{Latin Modern Math}
\usepackage{enumitem,needspace,xurl,hyperref}
\setlength{\parindent}{0pt}
\setlength{\parskip}{5pt}
\setlength{\emergencystretch}{3em}
\newcommand{\N}{\mathbb{N}}
\newcommand{\Z}{\mathbb{Z}}
\newcommand{\Q}{\mathbb{Q}}
\newcommand{\R}{\mathbb{R}}
\newcommand{\C}{\mathbb{C}}
\newcommand{\F}{\mathbb{F}}
\newcommand{\A}{\mathbb{A}}
\newcommand{\PP}{\mathbb{P}}
\newcommand{\E}{\mathbb{E}}
\newcommand{\eps}{\varepsilon}
\newcommand{\dd}{\,\mathrm{d}}
\begin{document}
\section*{UnsolvedMath ID: 801002}
\textbf{Source problem code:} catalogue entry
\subsection{Definitions and mathematical statement}
A standard Borel space is a measurable space isomorphic to a Borel subset of a Polish space with its induced Borel structure. For this formulation, a \emph{Borel field} means a field \(F\) whose underlying measurable space is standard Borel and whose addition, multiplication, additive inverse, and multiplicative inverse on \(F\setminus\{0\}\) are Borel maps. The supplied extract does not explicitly define ``Borel field''; this conventional interpretation requires confirmation against the cited source.

A \emph{Borel algebraically closed extension} of a Borel field \(F\) consists of a Borel field \(K\) and a Borel injective field homomorphism \(\iota:F\to K\), preserving \(1\), such that \(K\) is algebraically closed. Explicitly, for every integer \(n\geq 1\) and every \(a_0,\ldots,a_n\in K\) with \(a_n\neq 0\), there exists \(z\in K\) such that
\[
\sum_{j=0}^{n}a_jz^j=0.
\]
The definition does not require \(K\) to be algebraic over \(\iota(F)\).

\textbf{Question (Montalb\'an--Nies, cited Question 2.8).} Does every Borel field \(F\) admit a Borel algebraically closed extension? Equivalently, for every Borel field \(F\), do there exist a standard Borel field \(K\) and a Borel injective unital field homomorphism \(\iota:F\to K\) such that every nonconstant polynomial over \(K\) has a root in \(K\)? No restriction on characteristic or cardinality is stated in the supplied record. The registry labels the problem open; that label is supplied source data, not an independently verified status assessment.
\subsection{Short English statement}
Can every field with a standard Borel structure and Borel field operations be embedded by a Borel field homomorphism into an algebraically closed Borel field? The larger field need not be algebraic over the embedded field. The meaning of ``Borel field'' should be confirmed against the source.
\subsection{Sources}
\begin{itemize}
\item Supplied local UnsolvedMath registry record: permanent ID \(801002\), source code \texttt{catalogue entry}, titled ``Open Problems in Mathematical Logic --- Problem 14.'' The registry records Barmpalias, Gao et al. as proposers, the year 2026, and status ``open.'' These are supplied registry attributes.
\item Barmpalias, Gao et al., \emph{Open Problems in Mathematical Logic} (2026), Problem 14. Supplied source URL: \url{https://arxiv.org/abs/2608.26628}. The background records access on 2026-10-16 and extraction by \texttt{arxiv\_latex}; these are supplied provenance claims, not independently verified access or extraction.
\item The extracted statement attributes the question to Montalb\'an--Nies, with reference label \texttt{MontalbanNies2013}, Question 2.8. The supplied record provides no full bibliographic entry or separate URL for that reference.
\end{itemize}
\end{document}
